\par\Needspace{4\baselineskip}
\begin{center}
\SourceLeaf{030b}\pdfbookmark[2]{12.}{article-012}12.
\end{center}
Ein Näherungswert für $E$, mit dem man die Rechnung beginnen kann, wird sich meist leicht finden lassen, besonders wenn die Aufgabe für mehrere Werte von $M$ zu lösen und für einige davon bereits gelöst ist. Fehlen alle anderen Hilfsmittel, so steht wenigstens fest, dass $E$ zwischen den Grenzen $M$ und $M \pm e$ liegen muss (wobei die Exzentrizität $e$ in Bogensekunden ausgedrückt wird und das obere Vorzeichen im ersten und zweiten Quadranten, das untere im dritten und vierten gilt). Daher kann man als Anfangswert für $E$ entweder $M$ selbst oder einen nach einer groben Schätzung vergrößerten oder verkleinerten Wert ansetzen. Es versteht sich fast von selbst, dass die erste Rechnung bei einem wenig genauen Anfangswert keine besondere Genauigkeit verlangt und dass kleinere Tafeln, wie sie der berühmte \textsc{Lalande} erstellt hat, völlig ausreichen. Zur Erleichterung der Rechnung wird man ferner für $\epsilon$ stets Werte wählen, deren Sinus sich ohne Interpolation unmittelbar den Tafeln entnehmen lassen. Man kann etwa volle Bogenminuten oder Vielfache von zehn Bogensekunden wählen, je nachdem, ob die verwendeten Tafeln in Schritten von einer Bogenminute oder von zehn Bogensekunden fortschreiten. Wie diese Regeln anzupassen sind, wenn man die Winkel nach der neuen Dezimalteilung angibt, kann sich jeder selbst überlegen.
